\section{Open Questions}
\label{sec:open}

Five scientifically important questions that this replication does
\emph{not} settle. Each is grounded in a re-read of Turner et al.\ 2019
against the reproduced pipeline in this directory.

\begin{enumerate}
\item \textbf{Partial-body / non-uniform contamination geometry.}
Does the $(F,t)\!\to\!A$ inversion generalize when $^{137}$Cs is
non-uniformly distributed (e.g.\ sequestered in liver/kidney with
heterogeneous dose to circulating lymphocytes)?
\emph{Basis:} the paper's Eq.~3 assumes whole-body well-mixed
biokinetics from a single IV bolus; real accidental intakes
(inhalation, ingestion) give heterogeneous organ deposition and
blood/spleen may see very different local dose-rates than the
whole-body average. Nothing in the reproduced pipeline exercises
non-uniform geometry.
\emph{Next steps:} simulate a two-compartment biokinetic model
(GI + soft-tissue 30/70 split, or ICRP-137 inhalation type-M) and
re-run $F(A,t)$ inversion on synthetic blood; check whether the
calibration curve degrades and by how much AUC drops.

\item \textbf{Protracted-dose vs acute equivalence.}
Is external protracted-dose exposure at the same integrated dose
equivalent to the internal-emitter acute-then-decaying case that the
paper fits?
\emph{Basis:} $F(A,t) = b + kAt\exp(Q_1{+}Q_2)$ explicitly encodes
ongoing exposure via $kAt$, but external-photon biodosimetry
(LNT/DDREF) says protracted dose deposits less biological damage per
Gy than acute. The paper does not test the equivalence.
\emph{Next steps:} pair the internal cohort with an external Cs-137
protracted exposure at 0.5--5~Gy over 2--14~d and compare
$\gamma$-H2AX$(t)$ against the fit; test whether one $k$ suffices or
a scaling factor is needed.

\item \textbf{Inter-donor CV floor and true LoD.}
What is the per-mouse coefficient-of-variation of $\gamma$-H2AX total
fluorescence, and does it set a hard limit below which the
inversion cannot resolve nearby activities?
\emph{Basis:} group SEM ($n{=}8$) is reported but individual-mouse
values are not in the supplements. The reported ROC AUC $=0.93$ may
be inflated by group averaging; a real deployment sees one individual.
\emph{Next steps:} obtain per-mouse fluorescence values, OR run
$n{\ge}12$ mice at closely-spaced activities (6.66 vs 7.65 MBq),
measure CV directly, and recompute AUC with $n\!\approx\!160$ vs
$n\!=\!20$ group means.

\item \textbf{Integration with cytogenetic/PCC biodosimetry.}
Does combining $\gamma$-H2AX with dicentric-chromosome or PCC assays
improve the poor day-14 correlation (Pearson 0.34)?
\emph{Basis:} $\gamma$-H2AX resolves early (h--d) DSB damage;
PCC/dicentrics have different kinetics and different LoDs. The paper
positions $\gamma$-H2AX as standalone; a joint readout is untested.
\emph{Next steps:} paired-assay cohort ($\gamma$-H2AX + PCC on same
draw at days 2, 5, 14). Fit joint $(F_{\gamma},F_{\mathrm{PCC}},t)
\!\to\!A$ and test whether joint AUC exceeds the marginal 0.93.

\item \textbf{True low-dose LoD below 0.1~Gy.}
Can the pipeline discriminate against sham controls in the operationally
relevant 0.1--0.5~Gy region?
\emph{Basis:} the paper's lowest activity (5.74~MBq) accrues
$\sim$3.4~Gy to day 14 — an order of magnitude above the ICRP/WHO
triage cut-off. The paper does not present control-vs-lowest ROC.
\emph{Next steps:} extend the mouse cohort with three low-activity
groups (0.5, 1.0, 2.5~MBq) plus $n{\ge}12$ sham controls, sampled at
days 2--3 (best window). Report AUC vs control at each activity and
publish an explicit LoD.
\end{enumerate}
